This chapter consolidates every performance measurement reported in this
work into a single narrative, in the order a reviewer would ask about
them: \emph{what questions did we set out to answer, which experiments
answer each, and what did we find?} It is the IMRAD ``Results'' chapter
of an MDPI-style structured article and does not duplicate the subsystem
narratives in \url{memcpy.md} and \url{cam_switch.md}; those remain the
primary source for the implementation detail behind each result.

Subsection numbering in the captions of the tables below matches the
subsections of this chapter.

\begin{center}\rule{0.5\linewidth}{0.5pt}\end{center}

\section{1. Research questions}\label{research-questions}

We asked five questions of the platform, each addressable by one or more
experiments in the \texttt{diag} framework (see
\url{experimental_setup.md} §6.1).

{\def\LTcaptype{none} % do not increment counter
\begin{longtable}[]{@{}
  >{\raggedright\arraybackslash}p{(\linewidth - 4\tabcolsep) * \real{0.3333}}
  >{\raggedright\arraybackslash}p{(\linewidth - 4\tabcolsep) * \real{0.3333}}
  >{\raggedright\arraybackslash}p{(\linewidth - 4\tabcolsep) * \real{0.3333}}@{}}
\toprule\noalign{}
\begin{minipage}[b]{\linewidth}\raggedright
\#
\end{minipage} & \begin{minipage}[b]{\linewidth}\raggedright
Question
\end{minipage} & \begin{minipage}[b]{\linewidth}\raggedright
Experiments that answer it
\end{minipage} \\
\midrule\noalign{}
\endhead
\bottomrule\noalign{}
\endlastfoot
\textbf{RQ1} & What sustained frame rate does the parallel vision
pipeline achieve on a fixed camera, for the 512×512 single-class model?
& E14 (baseline 80-class model) → E15 (target model) \\
\textbf{RQ2} & Can the SDRAM-to-tensor copy --- identified as the
dominant stage of InferTask --- be accelerated without changing the
pipeline's producer-consumer contract? & E15 \texttt{dma\_memcpy(0/1)}
A/B, same firmware \\
\textbf{RQ3} & What is the cost of switching between the two cameras on
a frame-by-frame basis on a single-MIPI-lane-with-MUX topology? & E16
alternating; E18 three-sweep head-to-tail \\
\textbf{RQ4} & Can that switch be made \textbf{glitch-free} at the pixel
level without a hardware rework? & E17 per-switch JPEG visual
inspection \\
\textbf{RQ5} & Do the platform's fault-handling pathways leave
observable traces when a measurement runs without problems, and how do
they escalate when the system degrades? & fault counters exposed via
\texttt{sentai.diag.cam\_stats()}, observed at the end of every E18
session \\
\end{longtable}
}

The rest of this chapter presents, per question, the headline
measurement, the cross-session reproducibility check, and a brief
interpretation that points to the implementation chapter for depth.

\begin{center}\rule{0.5\linewidth}{0.5pt}\end{center}

\section{2. RQ1 --- Parallel-pipeline steady-state
throughput}\label{rq1-parallel-pipeline-steady-state-throughput}

\subsection{2.1 Setup}\label{setup}

Experiment class E15 (\texttt{e15\_pipeline\_parallel\_512}) drives
\texttt{sentai.pipeline.start/get/stop} with \texttt{PrepTask} and
\texttt{InferTask} running concurrently. \texttt{PrepTask} does
frame-grab + PXP resize + int8 quantisation into a staging buffer;
\texttt{InferTask} copies the staging buffer into the TFLite input
tensor, then calls Invoke and NMS. The two tasks communicate through the
\texttt{staging\_free} / \texttt{prep\_done} semaphore pair documented
in \url{camera.md} §PXP. The wall-clock interval between two successive
\texttt{sentai.pipeline.get()} returns is the per-frame throughput.

20 runs of 20 iterations each per session; all reported numbers are
post-eDMA (see RQ2 for the pre/post split).

\subsection{2.2 Headline result (Table
1)}\label{headline-result-table-1}

\textbf{Table 1.} Sustained parallel-pipeline throughput on the 512×512
single-class model, build \#622+ (post-eDMA memcpy). Source:
\href{../experiments/s032_e15_x20/}{experiments/s032\_e15\_x20/} --- 20
runs, 20 repetitions each.

{\def\LTcaptype{none} % do not increment counter
\begin{longtable}[]{@{}lrr@{}}
\toprule\noalign{}
Metric & Mean ± σ & min / max \\
\midrule\noalign{}
\endhead
\bottomrule\noalign{}
\endlastfoot
\texttt{frame\_interval\_ms} & 64.7 ± 1.9 & 61 / 75 \\
\texttt{prep\_stall\_ms} & 0.1 ± 0.3 & 0 / 1 \\
\texttt{infer\_stall\_ms} & 41.1 ± 1.5 & 37 / 45 \\
\textbf{Sustained FPS} & \textbf{15.46 ± 0.45} & --- \\
\end{longtable}
}

\subsection{2.3 Interpretation}\label{interpretation}

The pipeline is \textbf{stall-balanced}.
\texttt{infer\_stall\_ms\ ≈\ 41\ ms} means that InferTask blocks for
\textasciitilde41 ms per frame waiting for PrepTask to deliver the next
staging buffer; that is, PrepTask is the critical- path stage
(\textasciitilde64 ms/frame), and InferTask (Invoke + NMS + memcpy)
finishes in ≈ 23 ms and waits. The 15.5 FPS ceiling therefore reflects
the PrepTask stage time, not the TPU. The 2 FPS margin to the 15 FPS
sensor rate is explained by the sensor actually streaming at slightly
above 15 Hz (ISR count per second observed is 15.2-15.5 on this
hardware).

\begin{center}\rule{0.5\linewidth}{0.5pt}\end{center}

\section{3. RQ2 --- eDMA acceleration of the staging-to-tensor
copy}\label{rq2-edma-acceleration-of-the-staging-to-tensor-copy}

\subsection{3.1 Setup}\label{setup-1}

A runtime A/B flag
\texttt{sentai.pipeline.dma\_memcpy(0\ \textbar{}\ 1)} selects between
the baseline CPU \texttt{memcpy()} and the eDMA-based path (channel 31,
32-byte AXI bursts). Both settings coexist in the same firmware build;
switching between them is a single \texttt{volatile} store and does not
restart the pipeline. 10 × 20 frames per setting, captured in the same
session so scene, thermal state and TPU package cache are held constant.

\subsection{3.2 Headline result (Table
2)}\label{headline-result-table-2}

\textbf{Table 2.} Per-stage timing and sustained FPS for CPU
\texttt{memcpy} vs eDMA memcpy. 10 runs × 20 frames each, paired within
a single session. Source:
\href{../_e15_ab.py}{\texttt{../\_e15\_ab.py}}, detail tables in
\url{memcpy.md} §``Per-run results''.

{\def\LTcaptype{none} % do not increment counter
\begin{longtable}[]{@{}
  >{\raggedleft\arraybackslash}p{(\linewidth - 12\tabcolsep) * \real{0.1429}}
  >{\raggedleft\arraybackslash}p{(\linewidth - 12\tabcolsep) * \real{0.1429}}
  >{\raggedleft\arraybackslash}p{(\linewidth - 12\tabcolsep) * \real{0.1429}}
  >{\raggedleft\arraybackslash}p{(\linewidth - 12\tabcolsep) * \real{0.1429}}
  >{\raggedleft\arraybackslash}p{(\linewidth - 12\tabcolsep) * \real{0.1429}}
  >{\raggedleft\arraybackslash}p{(\linewidth - 12\tabcolsep) * \real{0.1429}}
  >{\raggedleft\arraybackslash}p{(\linewidth - 12\tabcolsep) * \real{0.1429}}@{}}
\toprule\noalign{}
\begin{minipage}[b]{\linewidth}\raggedleft
\texttt{dma\_memcpy}
\end{minipage} & \begin{minipage}[b]{\linewidth}\raggedleft
Invoke (ms)
\end{minipage} & \begin{minipage}[b]{\linewidth}\raggedleft
memcpy (ms)
\end{minipage} & \begin{minipage}[b]{\linewidth}\raggedleft
NMS (ms)
\end{minipage} & \begin{minipage}[b]{\linewidth}\raggedleft
total\_infer (ms)
\end{minipage} & \begin{minipage}[b]{\linewidth}\raggedleft
Wall (ms)
\end{minipage} & \begin{minipage}[b]{\linewidth}\raggedleft
\textbf{FPS}
\end{minipage} \\
\midrule\noalign{}
\endhead
\bottomrule\noalign{}
\endlastfoot
0 (baseline) & 50.2 & \textbf{24.1} & 0.2 & 74.5 & 74.7 &
\textbf{13.39} \\
1 (optimised) & 49.7 & \textbf{14.6} & 0.2 & 64.5 & 65.1 &
\textbf{15.32} \\
Δ & −0.5 & \textbf{−9.5} & 0 & \textbf{−10.0} & \textbf{−9.6} &
\textbf{+1.93} \\
\end{longtable}
}

\subsection{3.3 Interpretation}\label{interpretation-1}

The memcpy drops from 32 \% of the per-frame critical path to 22 \%.
Invoke time is unchanged (TPU is not involved in the memcpy), so the
improvement is wholly attributable to the eDMA. The 9.6 ms wall- time
saving translates 1-to-1 into FPS because we are on the critical path of
InferTask. The optimisation crossed the 15 FPS sensor-rate threshold,
which is why later camera-switch experiments can treat ``fixed camera,
30 fps sensor, \textasciitilde16 FPS pipeline output'' as a baseline.
Implementation detail --- eDMA configuration, cache-coherency notes, two
debugging trips --- is documented in full in \url{memcpy.md}
§``Optimised''.

\begin{center}\rule{0.5\linewidth}{0.5pt}\end{center}

\section{4. RQ3 --- Per-switch cost on the shared-MIPI, MUX-gated dual
sensor}\label{rq3-per-switch-cost-on-the-shared-mipi-mux-gated-dual-sensor}

\subsection{4.1 Setup}\label{setup-2}

Experiment E18 (\texttt{e18\_camera\_switch\_headtail}) runs three
sweeps of the same sequential loop
(\texttt{select\ →\ to\_tensor\ →\ invoke\ →\ detect}) back to back in
one session:

\begin{itemize}
\tightlist
\item
  \textbf{A} --- fixed on \texttt{cam\_a} for all 40 iterations (no
  switches)
\item
  \textbf{B} --- fixed on \texttt{cam\_b} for all 40 iterations (no
  switches)
\item
  \textbf{C} --- alternating \texttt{cam\_a\ ↔\ cam\_b} every iteration
\end{itemize}

Three sweeps let us express the switch cost as a subtraction against a
baseline measured under identical conditions:

\begin{verbatim}
per_switch_overhead = C_total_mean − max(A_total_mean, B_total_mean)
\end{verbatim}

\subsection{4.2 Headline result (Table
3)}\label{headline-result-table-3}

\textbf{Table 3.} Three-sweep head-to-tail benchmark, 30 fps sensor, Fix
B firmware (post-flip-on-EOF), \texttt{switch\_drain=2},
\texttt{ratio=(0,0)}, 40 iterations per sweep, first sample dropped.
Source:
\href{../experiments/s045_e18_post_refactor/}{experiments/s045\_e18\_post\_refactor/}.

{\def\LTcaptype{none} % do not increment counter
\begin{longtable}[]{@{}
  >{\raggedright\arraybackslash}p{(\linewidth - 12\tabcolsep) * \real{0.1111}}
  >{\raggedleft\arraybackslash}p{(\linewidth - 12\tabcolsep) * \real{0.1481}}
  >{\raggedleft\arraybackslash}p{(\linewidth - 12\tabcolsep) * \real{0.1481}}
  >{\raggedleft\arraybackslash}p{(\linewidth - 12\tabcolsep) * \real{0.1481}}
  >{\raggedleft\arraybackslash}p{(\linewidth - 12\tabcolsep) * \real{0.1481}}
  >{\raggedleft\arraybackslash}p{(\linewidth - 12\tabcolsep) * \real{0.1481}}
  >{\raggedleft\arraybackslash}p{(\linewidth - 12\tabcolsep) * \real{0.1481}}@{}}
\toprule\noalign{}
\begin{minipage}[b]{\linewidth}\raggedright
Sweep
\end{minipage} & \begin{minipage}[b]{\linewidth}\raggedleft
\texttt{select}
\end{minipage} & \begin{minipage}[b]{\linewidth}\raggedleft
\texttt{to\_tensor}
\end{minipage} & \begin{minipage}[b]{\linewidth}\raggedleft
\texttt{invoke}
\end{minipage} & \begin{minipage}[b]{\linewidth}\raggedleft
\texttt{detect}
\end{minipage} & \begin{minipage}[b]{\linewidth}\raggedleft
\textbf{total (ms)}
\end{minipage} & \begin{minipage}[b]{\linewidth}\raggedleft
\textbf{FPS}
\end{minipage} \\
\midrule\noalign{}
\endhead
\bottomrule\noalign{}
\endlastfoot
A --- fixed cam0 & 0.0 & 31.6 & 29.8 & 0.4 & \textbf{61.8} &
\textbf{16.18} \\
B --- fixed cam1 & 0.0 & 31.6 & 30.5 & 0.4 & \textbf{62.5} &
\textbf{16.00} \\
C --- alternating & 17.6 & 96.0 & 31.5 & 0.7 & \textbf{145.8} &
\textbf{6.86} \\
\end{longtable}
}

\textbf{Per-switch overhead = 83.3 ms} (133.4 \% of the slower
baseline). \textbf{Directional asymmetry (cam1 − cam0) = −2.1 ms}
(within noise).

\subsection{4.3 Budget decomposition}\label{budget-decomposition}

{\def\LTcaptype{none} % do not increment counter
\begin{longtable}[]{@{}
  >{\raggedright\arraybackslash}p{(\linewidth - 2\tabcolsep) * \real{0.4286}}
  >{\raggedleft\arraybackslash}p{(\linewidth - 2\tabcolsep) * \real{0.5714}}@{}}
\toprule\noalign{}
\begin{minipage}[b]{\linewidth}\raggedright
Component of the 83.3 ms tax
\end{minipage} & \begin{minipage}[b]{\linewidth}\raggedleft
Amount
\end{minipage} \\
\midrule\noalign{}
\endhead
\bottomrule\noalign{}
\endlastfoot
\texttt{select} wait for EOF ISR consumption & 17.6 ms (≈ 0.5 frame
interval at 30 fps) \\
\texttt{to\_tensor} post-switch drain (stale-buffer drop + 2 fresh
frames) & 64.4 ms (≈ 2 × 33 ms) \\
Sensor-independent stages (\texttt{invoke}, \texttt{detect}) & 0 ms
(same in A and C) \\
\end{longtable}
}

\subsection{4.4 Cross-session agreement}\label{cross-session-agreement}

Three sessions captured the same three-sweep benchmark under
successively more refined firmware (pre-refactor, post-refactor
confirmation, NASA-JPL review fixes). Agreement across the three
sessions is a reproducibility check (Table 4).

\textbf{Table 4.} Cross-session reproducibility of Table 3. All three
rows are 40-iteration head-to-tail sweeps on the same scene/model;
firmware differences are the only variable. Values in ms.

{\def\LTcaptype{none} % do not increment counter
\begin{longtable}[]{@{}
  >{\raggedright\arraybackslash}p{(\linewidth - 8\tabcolsep) * \real{0.1579}}
  >{\raggedleft\arraybackslash}p{(\linewidth - 8\tabcolsep) * \real{0.2105}}
  >{\raggedleft\arraybackslash}p{(\linewidth - 8\tabcolsep) * \real{0.2105}}
  >{\raggedleft\arraybackslash}p{(\linewidth - 8\tabcolsep) * \real{0.2105}}
  >{\raggedleft\arraybackslash}p{(\linewidth - 8\tabcolsep) * \real{0.2105}}@{}}
\toprule\noalign{}
\begin{minipage}[b]{\linewidth}\raggedright
Session
\end{minipage} & \begin{minipage}[b]{\linewidth}\raggedleft
A --- fixed cam0
\end{minipage} & \begin{minipage}[b]{\linewidth}\raggedleft
B --- fixed cam1
\end{minipage} & \begin{minipage}[b]{\linewidth}\raggedleft
C --- alternating
\end{minipage} & \begin{minipage}[b]{\linewidth}\raggedleft
Asymmetry cam1 − cam0
\end{minipage} \\
\midrule\noalign{}
\endhead
\bottomrule\noalign{}
\endlastfoot
\href{../experiments/s043_e18_headtail_drain2/}{\texttt{s043\_e18\_headtail\_drain2}}
& 62.7 & 62.5 & 145.3 & +0.1 \\
\href{../experiments/s044_e18_headtail_drain2/}{\texttt{s044\_e18\_headtail\_drain2}}
& 62.6 & 62.1 & 145.7 & +0.1 \\
\href{../experiments/s045_e18_post_refactor/}{\texttt{s045\_e18\_post\_refactor}}
& 61.8 & 62.5 & 145.8 & −2.1 \\
\textbf{σ across sessions} & \textbf{0.4 ms} & \textbf{0.2 ms} &
\textbf{0.3 ms} & noise \\
\end{longtable}
}

The spread across three firmware builds is \textless{} 1 ms on every
stage. That bounds the measurement noise and validates the claim that
the NASA-JPL review fixes are performance-neutral.

\begin{center}\rule{0.5\linewidth}{0.5pt}\end{center}

\section{5. RQ4 --- Glitch-free MUX
transitions}\label{rq4-glitch-free-mux-transitions}

\subsection{5.1 Setup}\label{setup-3}

E17 (\texttt{e17\_switch\_drain\_visual}) writes alternating-camera
JPEGs into the MicroPython heap during the timing loop and flushes them
to LittleFS afterwards. Two runs per session at \texttt{switch\_drain=2}
and \texttt{switch\_drain=1}, 16 iterations each, 512×512 quality-70.
The output is a set of pixel-level frames available for offline review.

The subject is \emph{visual correctness}, not timing. The timing is a
by-product.

\subsection{5.2 Headline result --- visual
inspection}\label{headline-result-visual-inspection}

\textbf{Table 5.} Visual inspection outcomes across firmware builds. The
full frame sets are in the linked folders; a representative ``broken''
frame and its ``fixed'' counterpart are cited inline.

{\def\LTcaptype{none} % do not increment counter
\begin{longtable}[]{@{}
  >{\raggedright\arraybackslash}p{(\linewidth - 6\tabcolsep) * \real{0.2500}}
  >{\raggedright\arraybackslash}p{(\linewidth - 6\tabcolsep) * \real{0.2500}}
  >{\raggedright\arraybackslash}p{(\linewidth - 6\tabcolsep) * \real{0.2500}}
  >{\raggedright\arraybackslash}p{(\linewidth - 6\tabcolsep) * \real{0.2500}}@{}}
\toprule\noalign{}
\begin{minipage}[b]{\linewidth}\raggedright
Firmware
\end{minipage} & \begin{minipage}[b]{\linewidth}\raggedright
Session
\end{minipage} & \begin{minipage}[b]{\linewidth}\raggedright
drain=1 outcome
\end{minipage} & \begin{minipage}[b]{\linewidth}\raggedright
drain=2 outcome
\end{minipage} \\
\midrule\noalign{}
\endhead
\bottomrule\noalign{}
\endlastfoot
Pre-flip-on-EOF &
\href{../experiments/s038_e17_drain_ab/}{\texttt{s038\_e17\_drain\_ab}}
& \textbf{Fails.} Horizontal seam at 40-60 \% image height, every other
frame split between cam0 and cam1. Example:
\href{../experiments/s038_e17_drain_ab/e17_t1_frames/002_cam0_133ms.jpg}{\texttt{002\_cam0\_133ms.jpg}}
& Clean. \\
Post-flip-on-EOF &
\href{../experiments/s041_e17_eof_check/}{\texttt{s041\_e17\_eof\_check}}
& \textbf{Mid-buffer seam removed}, but residual sensor-side artefacts
remain (AEC/AGC convergence) --- see \url{threats_to_validity.md} §3. &
Clean. Example:
\href{../experiments/s041_e17_eof_check/e17_t1_frames/003_cam1_201ms.jpg}{\texttt{003\_cam1\_201ms.jpg}} \\
\end{longtable}
}

\subsection{5.3 Interpretation}\label{interpretation-2}

Moving the GPIO flip into the CSI end-of-frame ISR --- so the MUX
transition lands in the MIPI VBLANK window between two DMA buffers ---
eliminates the class of failure in which a single DMA buffer contains
pixels from two sensors. The shipping configuration is
\texttt{switch\_drain=2}: it produces reliably clean frames.
\texttt{switch\_drain=1} remains available as an experimentation hook
but the frames it produces are not suitable for production use; this is
the only residual pixel-quality limitation of the camera-switch chapter
and is called out in its own ``Known limitations'' section in
\url{cam_switch.md}.

\begin{center}\rule{0.5\linewidth}{0.5pt}\end{center}

\section{6. RQ5 --- Fault observability}\label{rq5-fault-observability}

\subsection{6.1 Setup}\label{setup-4}

Every degraded path in the camera-switch subsystem --- ISR arm not
consumed within 150 ms, drain wait hitting its 300 ms ceiling,
\texttt{GetRawFrame} retry, \texttt{GetRawFrame} fatal --- increments a
persistent counter (see the \texttt{0x0Axx} range in
\href{../error_codes.csv}{error\_codes.csv}). The counters are exposed
to MicroPython as \texttt{sentai.diag.cam\_stats()}. Every E18 session
captures the counter dict at the end of the run; a non-zero
degraded-path value invalidates that session's interpretation.

\subsection{6.2 Headline result}\label{headline-result}

\textbf{Table 6.} Post-run fault counters for the canonical E18 sessions
cited in this paper. All zeros means every switch was handled by the
nominal flip-on-EOF path.

{\def\LTcaptype{none} % do not increment counter
\begin{longtable}[]{@{}
  >{\raggedright\arraybackslash}p{(\linewidth - 10\tabcolsep) * \real{0.1304}}
  >{\raggedleft\arraybackslash}p{(\linewidth - 10\tabcolsep) * \real{0.1739}}
  >{\raggedleft\arraybackslash}p{(\linewidth - 10\tabcolsep) * \real{0.1739}}
  >{\raggedleft\arraybackslash}p{(\linewidth - 10\tabcolsep) * \real{0.1739}}
  >{\raggedleft\arraybackslash}p{(\linewidth - 10\tabcolsep) * \real{0.1739}}
  >{\raggedleft\arraybackslash}p{(\linewidth - 10\tabcolsep) * \real{0.1739}}@{}}
\toprule\noalign{}
\begin{minipage}[b]{\linewidth}\raggedright
Session
\end{minipage} & \begin{minipage}[b]{\linewidth}\raggedleft
\texttt{switch\_ok\_eof}
\end{minipage} & \begin{minipage}[b]{\linewidth}\raggedleft
\texttt{switch\_fallback}
\end{minipage} & \begin{minipage}[b]{\linewidth}\raggedleft
\texttt{drain\_timeout}
\end{minipage} & \begin{minipage}[b]{\linewidth}\raggedleft
\texttt{grab\_retry}
\end{minipage} & \begin{minipage}[b]{\linewidth}\raggedleft
\texttt{grab\_fatal}
\end{minipage} \\
\midrule\noalign{}
\endhead
\bottomrule\noalign{}
\endlastfoot
\texttt{s043\_e18\_headtail\_drain2} & 41 & \textbf{0} & \textbf{0} &
\textbf{0} & \textbf{0} \\
\texttt{s044\_e18\_headtail\_drain2} & 41 & \textbf{0} & \textbf{0} &
\textbf{0} & \textbf{0} \\
\texttt{s045\_e18\_post\_refactor} & 41 & \textbf{0} & \textbf{0} &
\textbf{0} & \textbf{0} \\
\end{longtable}
}

\subsection{6.3 Interpretation}\label{interpretation-3}

The fault-counter infrastructure is not merely defensive
instrumentation; it is a publication-grade \textbf{measurement integrity
check}. A reviewer who suspects that some fraction of our E18
alternating iterations fell back to the legacy synchronous switch path
(which would re-introduce the seam and bias the timing) can confirm from
the counters that in the sessions we report, \textbf{none did}. The
counters remain available at all times via REPL for future replays on
the same hardware.

\begin{center}\rule{0.5\linewidth}{0.5pt}\end{center}

\section{7. Summary table (Research-question ×
metric)}\label{summary-table-research-question-metric}

\textbf{Table 7.} Paper-level summary: one row per research question,
the experiment class that answers it, the headline number, the direction
of change versus the baseline, and the session(s) that document it.

{\def\LTcaptype{none} % do not increment counter
\begin{longtable}[]{@{}
  >{\raggedright\arraybackslash}p{(\linewidth - 10\tabcolsep) * \real{0.1429}}
  >{\raggedright\arraybackslash}p{(\linewidth - 10\tabcolsep) * \real{0.1429}}
  >{\raggedleft\arraybackslash}p{(\linewidth - 10\tabcolsep) * \real{0.1905}}
  >{\raggedleft\arraybackslash}p{(\linewidth - 10\tabcolsep) * \real{0.1905}}
  >{\raggedleft\arraybackslash}p{(\linewidth - 10\tabcolsep) * \real{0.1905}}
  >{\raggedright\arraybackslash}p{(\linewidth - 10\tabcolsep) * \real{0.1429}}@{}}
\toprule\noalign{}
\begin{minipage}[b]{\linewidth}\raggedright
RQ
\end{minipage} & \begin{minipage}[b]{\linewidth}\raggedright
Headline metric
\end{minipage} & \begin{minipage}[b]{\linewidth}\raggedleft
Baseline
\end{minipage} & \begin{minipage}[b]{\linewidth}\raggedleft
Final
\end{minipage} & \begin{minipage}[b]{\linewidth}\raggedleft
Δ
\end{minipage} & \begin{minipage}[b]{\linewidth}\raggedright
Evidence
\end{minipage} \\
\midrule\noalign{}
\endhead
\bottomrule\noalign{}
\endlastfoot
RQ1 & Sustained parallel-pipeline FPS, 512×512 & 13.39 (pre-eDMA) &
15.47 & \textbf{+15 \%} & \href{../experiments/s024_e15_x20/}{s024},
\href{../experiments/s026_e15_x20/}{s026},
\href{../experiments/s032_e15_x20/}{s032} \\
RQ2 & Staging→tensor memcpy cost & 24.1 ms (CPU) & 14.6 ms (eDMA) &
\textbf{−39 \%} & \href{memcpy.md}{memcpy.md §``Per-run results''} \\
RQ3 & Alternating FPS (switch every frame) & 4.7 (15 fps, pre-Fix A) &
6.86 (30 fps, Fix B) & \textbf{+46 \%} &
\href{../experiments/s034_e15_vs_e16_x40/}{s034},
\href{../experiments/s045_e18_post_refactor/}{s045} \\
RQ3 & Directional asymmetry cam1 − cam0 & +64.2 ms & ≤ 2.1 ms (within
noise) & \textbf{removed} &
\href{../experiments/s034_e15_vs_e16_x40/}{s034} →
\href{../experiments/s045_e18_post_refactor/}{s045} \\
RQ4 & Mid-buffer seam at drain=2 & visible every switch & none &
\textbf{removed} & visual inspection of
\href{../experiments/s038_e17_drain_ab/}{s038} vs
\href{../experiments/s041_e17_eof_check/}{s041} \\
RQ5 & Degraded-path counter during a nominal run & n/a pre-review & 0
across all three reported sessions & \textbf{validated} &
\url{artifact.md} §``Data integrity'' \\
\end{longtable}
}

\begin{center}\rule{0.5\linewidth}{0.5pt}\end{center}

\section{\texorpdfstring{8. Discussion (of results only; full discussion
in
\url{discussion.md})}{8. Discussion (of results only; full discussion in discussion.md)}}\label{discussion-of-results-only-full-discussion-in-discussion.md}

Three observations that belong here rather than in the implementation
chapters because they relate to the results taken together:

\begin{enumerate}
\def\labelenumi{\arabic{enumi}.}
\tightlist
\item
  \textbf{Each optimisation exposes the next bottleneck.} The CPU memcpy
  was hidden behind a slow camera-drain fix (described in
  \url{camera.md}). The 15-fps pipeline plateau was hidden behind the
  memcpy. The alternating-switch tax was hidden behind the fixed-camera
  ceiling. Each of the three ``headline'' results only became measurable
  once the previous layer had stopped dominating the wall-clock; this is
  why the subsystem narratives must be read in order (camera-drain →
  memcpy → cam\_switch).
\item
  \textbf{The residual 83 ms per-switch tax is sensor-topology bound,
  not firmware bound.} With both cameras sharing one CSI-2 lane, two
  frame intervals of drain is the theoretical floor for a conservative
  drain threshold on a not-VSYNC-synchronised pair of sensors. Further
  reduction requires either hardware sync (a FSIN wire between the
  OV5640s, rejected for this iteration --- see
  \url{threats_to_validity.md} §4) or a different CSI-receiver topology.
\item
  \textbf{The measurement discipline scales.} The same \texttt{diag}
  framework that captured the initial 6 FPS baseline in the E10-E12
  subsystem probes (Appendix A of
  \href{../experiments/README.md}{experiments/README.md}) captured the
  final 16 FPS head-to-tail benchmark. No measurement methodology change
  was needed to follow the four-ish orders of magnitude of scope from
  ``does \texttt{pvPortMalloc} work under load'' to ``what is the
  per-switch overhead of the dual-sensor MUX topology''. That is a
  statement about the framework, not about the results.
\end{enumerate}

Open questions, limitations, and threats to this interpretation are in
\url{threats_to_validity.md}.

\begin{center}\rule{0.5\linewidth}{0.5pt}\end{center}

\section{Data availability statement}\label{data-availability-statement}

Every number in Tables 1--7 is derived from raw per-iteration CSVs that
are mirrored under \href{../experiments/}{\texttt{../experiments/}}. The
session names cited in the table captions are the direct folder paths.
The appendix generator at
\href{../experiments/_build_appendix.py}{\texttt{../experiments/\_build\_appendix.py}}
reproduces the full appendix tables from those CSVs without any
on-device round-trip, so an external reviewer can verify every derived
statistic offline. The full artifact description --- code commit SHAs,
firmware build numbers, reproduction steps --- is in \url{artifact.md}.
